\section{The linear block transformation}
\label{free:section}

Exponential localization of Gaussian block kernels is a standard
ingredient of constructive renormalization. See
\cite{DybalskiStottmeisterTanimotoGreen} for a related treatment using
Fourier analyticity and resolvent estimates. We prove the particular
identities and the uniformity in the blocking factor required here.

The input is a Gaussian precision and a block average observed with
independent Gaussian noise. We first identify the new precision and the
conditional mean. We then express the same energy identity through
localized edge operators; this form will retain a positive cost near
large gradients in the nonlinear integration. Constants used to choose
the block size must be uniform in that size. The Fourier estimates in
Appendix~\ref{app:analytic-block} establish this uniformity on fixed
complex domains, so their use below does not shrink a domain at each
iteration.

Fix a sufficiently large dyadic integer $L$. Within each $L\times L$
block choose nonnegative weights $b_x$ obtained by sampling and
normalizing the product of the same even smooth one-dimensional bump.
If $x_{\rm c}$ is the center of the block, then
$\sum_xb_x(x-x_{\rm c})=0$ exactly. The bump is centered in the
block, is supported at distance at most $L/5$ from its center, and is
positive within distance $L/10$. Set
\[
 (Qu)(y)=\sum_{0\le x_1,x_2<L}b_xu(Ly+x),\qquad
 b(k)=\sum_xb_xe^{ik\cdot x}.
\]
The origins for the fine and coarse lattices are their block centers.
This convention removes an inessential translation in the affine
identities below. The observation precision is $a=1$; we retain $a$ in
formulas which are valid for any fixed positive value.

For Fourier momentum $k$, write
\[
 d_\mu(k)=e^{ik_\mu}-1,\qquad
 D(k)=\sum_{\mu=1}^2|d_\mu(k)|^2.
\]
Stars in formulas continued to complex momentum mean analytic
continuation of the adjoint on real momentum. At coarse momentum $p$ the
fine aliases are $k_l=(p+2\pi l)/L$, with $l$ a centered representative
modulo $L$ in each coordinate. Our Fourier convention assigns a fine
field the norm $L^2\sum_l|u_l|^2$ per coarse cell. Thus if a block-average
row has entries $b(k_l)$, its adjoint has entries $L^{-2}b(k_l)^*$.

\subsection{The free precision and the conditional mean}

Let $h$ count block transformations from the nearest-neighbor precision.
Define $K_0=D$ and
\begin{equation}
\label{free:recursion}
 K_{h+1}(p)^{-1}=a^{-1}
       +L^{-2}\sum_l|b(k_l)|^2K_h(k_l)^{-1}.
\end{equation}
The pole at zero in this expression is interpreted by continuation.
Equation~\eqref{free:recursion} is the covariance of $Qu$ plus independent
Gaussian observation noise.

\begin{lemma}
\label{free:bounds}
For all sufficiently large $L$, the following bounds hold with constants
independent of $h$ and $L$. The functions $K_h$ are analytic and uniformly
bounded on a common complex neighborhood of the momentum torus,
\[
 cD(p)\le K_h(p)\le CD(p)\quad(p\text{ real}),\qquad
 K_h(p)-D(p)=O(|p|^4).
\]
Their differences between depths $\bar h\ge h$ have the corresponding
analytic bounds multiplied by $L^{-h}$. There are real-kernel,
self-adjoint edge operators $B_h$, with analytic Hermitian Fourier
symbols on real momentum, such that
\[
 K_h=d^*B_hd,\qquad cI\le B_h\le CI,\qquad B_h(0)=I,
\]
whose inverses and history differences satisfy the same bounds.

Writing $K=K_h$ and $K'=K_{h+1}$, the conditional mean operator is
\begin{equation}
\label{free:P}
 P_l=L^{-2}K(k_l)^{-1}b(k_l)^*K'(p).
\end{equation}
It obeys $KP=Q^*K'$ and $a(I-QP)=K'$, preserves constants and centered
affine fields, and, for every fixed integer $r\ge0$,
\begin{equation}
\label{free:derivative}
 \sup_x\sum_y e^{c|x/L-y|}|\nabla_f^rP(x,y)|\le C_rL^{-r}.
\end{equation}
Here $\nabla_f$ is a nearest-neighbor difference on the fine lattice.
The conditional covariance
\begin{equation}
 C_h=(K_h+aQ^*Q)^{-1}
 \label{free:C}
\end{equation}
has exponentially decaying
kernels in coarse units, with constants allowed to depend on the fixed
$L$. All these assertions retain the factor $L^{-h}$ in history
differences.
\end{lemma}

\begin{proof}
The covariance iteration has an explicit form. Set $m=L^h$,
$W_h(\xi)=\prod_{j=1}^hb(\xi/L^j)$,
$D_m(p)=m^2D(p/m)$, and $s_L=\sum_xb_x^2$. Separating the principal
alias gives
\begin{equation}
\label{free:explicit}
 K_h(p)=\frac{D_m(p)}{W_h(p)W_h(p)^*+D_m(p)R_h(p)},
\end{equation}
where
\[
 R_h(p)=a^{-1}\sum_{r<h}s_L^r
 +\sum_{\substack{l\bmod m\\l\ne0}}
      \frac{W_h(p+2\pi l)W_h(p+2\pi l)^*}{D_m(p+2\pi l)}.
\]
The first sum is the accumulated observation noise: averaging $r$ more
times multiplies its variance by $s_L^r$. For the second sum,
Lemma~\ref{app:free-strip} proves that the product weights in the $r$th
alias shell have total at most $CA^r$, while its nonprincipal denominator
is at least $cL^{2(r-1)}$. Thus $R_h$ and its fixed analytic derivatives
are uniformly bounded. The principal bump product is uniformly positive
on the real momentum cube. These facts keep the denominator in
\eqref{free:explicit} away from zero on a common smaller complex domain.
They also give $K_h-D=O(|p|^4)$. Comparing the same products at two depths,
without dividing by possibly vanishing bump factors, gives the
$L^{-h}$ discrepancy and its fixed derivatives.

For the edge representation, the same lemma writes the fourth-order
zero $K_h-D$ as $d^*(B_h^{\rm raw}-I)d$, with
$B_h^{\rm raw}=I+O(|p|^2)$ and bounded analytic coefficients.
Hermitian, reality and spatial-symmetry averaging preserve this identity.
Adding a fixed sufficiently large multiple of $\mathsf c^*\mathsf c$, where
\[
 \mathsf c=(-d_2,d_1),\qquad \mathsf c d=0,
\]
does not change $d^*Bd$. Near zero the raw lift is already positive;
away from zero, ellipticity controls its longitudinal entry and the
added term controls the transverse entry. The fixed choice makes $B_h$
uniformly positive, with analytic inverse and the same history bounds.

The normal equations for the conditional mean give \eqref{free:P} and
the two stated operator identities. On a nonprincipal alias the smooth
bump gives arbitrary fixed inverse-power decay in $l$; for the principal
alias remove the apparent pole using
\[
 P_0=b(k_0)^{-1}\left[1-
 \left(a^{-1}+L^{-2}\sum_{l\ne0}|b(k_l)|^2K(k_l)^{-1}\right)K'\right].
\]
At $p=0$ the other aliases vanish to second order. The zero and first
jets of this formula prove preservation of constants and affine fields.
Each fine difference contributes $|d_\mu(k_l)|\le C(1+|l|)/L$.
Choose the bump decay exponent larger than the required number of
differences plus the summation dimension. Fourier inversion on a
smaller coarse strip proves \eqref{free:derivative}, with its weighted
moments and history version. This argument also treats differences
across block boundaries, since their coarse translation factor is
already in the alias formula.

Finally invert $K+aQ^*Q$ in aliases. The nonprincipal diagonal block is
invertible with the preceding bounds. Its rank-one update has
denominator
$a^{-1}+L^{-2}\sum_{l\ne0}|b(k_l)|^2K(k_l)^{-1}$, positive on real
momentum and nonzero on a smaller strip. The remaining principal Schur
complement, multiplied by $L^2$, is bounded below by the principal bump
term. Its inverse is analytic on that strip. Fourier inversion proves
the asserted localization of the covariance.
\end{proof}

The conditional mean also gives an exact energy decomposition. For real
ambient fields $u,V$ on a finite periodic pair of lattices,
\begin{equation}
 \frac12\langle u,K_hu\rangle+\frac a2\|V-Qu\|^2
 =\frac12\langle V,K_{h+1}V\rangle
  +\frac12\langle u-P_hV,C_h^{-1}(u-P_hV)\rangle.
 \label{free:schur-energy}
\end{equation}
Indeed $C_h^{-1}P_h=aQ^*$ and $a(I-QP_h)=K_{h+1}$ make the
cross and constant terms agree after expansion. Although $K_h$ vanishes
on constants, $C_h^{-1}=K_h+aQ^*Q$ is strictly positive, since $Q$
preserves constants. This identity identifies
the fluctuation field and its positive precision. The next construction
expresses the energy through localized edge rows, with the row and
column estimates needed to keep that positivity near rough regions.

\subsection{A positive identity for the kinetic energy}

The edge representation of $K_h$ is not unique: a curl correction leaves
$d^*B_hd$ unchanged. We use this freedom to obtain a positive completion
whose fine-output rows gain $L^{-1}$. All operators act on real ambient
fields, so the resulting identity also applies to sphere-valued fields
without linearizing them.

\begin{lemma}[Positive completion of the block transformation]
\label{lem:free-stack}
One can choose a common $c_0>0$ and edge operators $\ell_h$ such that
\[
 \mathsf A_h=\binom{\sqrt{c_0}I}{\ell_h},\qquad X_h=\mathsf A_hd,
\]
give the free precision $X_h^*X_h=K_h$. There are real operators
$\mathcal M_h,\mathcal N_h$ satisfying
\begin{align}
 \mathcal M_h^*\mathcal M_h+\mathcal N_h^*\mathcal N_h&=I,
 \label{free:isometry}\\
 \mathcal M_h^*(X_hu,\sqrt a(V-Qu))&=X_{h+1}V,
 \label{free:ambient}\\
 \mathcal M_hX_{h+1}&=(X_hP,\sqrt a(I-QP)),
 \label{free:harmonic}\\
 \mathcal N_hX_{h+1}&=0.
 \label{free:null}
\end{align}
All kernels have exponential moments and the history bounds of
Lemma~\ref{free:bounds}. The fine-output part of $\mathcal M_h$ has
weighted row norm $C/L$ and column norm $CL$. Its first two fine
differences have row norms $C/L^2,C/L^3$. All coarse slots have bounded
row and column norms. At depth zero one may take
$\ell_0=\sqrt{1-c_0}I$.
\end{lemma}

\begin{proof}
We first construct a coarse-edge to fine-edge map $T$ with
\begin{equation}
\label{free:Tidentities}
 T^*d=d'Q,\qquad TB'd'=BdP.
\end{equation}
Primes denote the next lattice. Start with
\[
 T_{0,l}=L^{-2}\operatorname{diag}(\omega_{\mu,l}^*)b(k_l)^*,\qquad
 \omega_{\mu,l}=\sum_{j=0}^{L-1}e^{ijk_{l,\mu}}.
\]
The identity $\omega_{\mu,l}d_\mu(k_l)=d'_\mu(p)$ proves the first
equation, with the $L^2$ convention for adjoints. Put
$R_l=B_ld_lP_l-T_{0,l}B'd'$. Its divergence vanishes by
$KP=Q^*K'$. Its aliases have bounds
$C_sL^{-1}(1+|l|)^{-s}$ for every required fixed $s$. At $p=0$ both
terms vanish on each nonprincipal alias; on the principal alias their
linear terms are both $ip/L$. Thus $R_0$ vanishes through first order.
Lemma~\ref{app:alias-division} applies and constructs analytic
$\gamma_l,H_l$ with
\[
 R_l=\mathsf c_l^*\gamma_l,\qquad H_ld'=\gamma_l,\qquad
 |\gamma_l|+|H_l|\le C_s(1+|l|)^{-s}.
\]
The first division removes a fine derivative. The second uses the
vanishing of $\gamma_l(0)$ on every alias to divide by coarse
derivatives, preserving alias gluing and decay without another power
of $L$. The interpolation formula and its uniform bound are proved in
Appendix~\ref{app:analytic-block}. Set
\[
 T_l=T_{0,l}+\mathsf c_l^*H_l(B')^{-1}.
\]
This proves \eqref{free:Tidentities}. Its aliases are bounded by
$C_JL^{-1}(1+|l|)^{-J}$. Fourier inversion gives the row, column, and
fine-difference bounds in the statement. The operations are linear or
uniformly analytic in the input coefficients, so also preserve the
history discrepancy. Symmetry averaging preserves both identities.

On real momentum the Hermitian defect
\[
 U=(B')^{-1}-T^*B^{-1}T-a^{-1}d'd'^*
\]
annihilates $B'd'$, by \eqref{free:Tidentities} and $a(I-QP)=K'$.
Division on both sides therefore factors $B'UB'$ through
$\mathsf c'^*\mathsf c'$ with a bounded analytic scalar coefficient.
Similarly $\mathsf c B^{-1}TB'd'=0$. Its aliases, before the last
division, cost $CL^{-2}(1+|l|)^{-J}$. Squaring and summing with the fine
norm $L^2\sum_l$ proves
\begin{align*}
 |z^*Uz|&\le C|\mathsf c'(B')^{-1}z|^2,\\
 \|\mathsf c B^{-1}Tz\|^2&\le CL^{-2}|\mathsf c'(B')^{-1}z|^2.
\end{align*}
The normalization factor in the second estimate is essential.

Choose a fixed sufficiently large $\lambda$ and replace
\[
 \widetilde B^{-1}=B^{-1}
     +\lambda B^{-1}\mathsf c^*\mathsf c B^{-1},
\]
and do the same at the next scale. This leaves
$\widetilde Bd=Bd$, hence leaves $K$ unchanged. The new defect
$U_*$ is bounded below by
\[
 [\lambda(1-C/L^2)-C]
       (B')^{-1}\mathsf c'^*\mathsf c'(B')^{-1}.
\]
First choose $\lambda\ge2(C+1)$ and then $L$ sufficiently large; the
coefficient is at least one. We may consequently write
\[
 \widetilde B'U_*\widetilde B'=\mathsf c'^*B_2\mathsf c'
\]
with a uniformly positive analytic scalar $B_2$. Positivity extends
through $p=0$, as follows from
\[
 \mathsf c'(B')^{-1}\widetilde B'
 =[1+\lambda\mathsf c'(B')^{-1}\mathsf c'^*]^{-1}\mathsf c'.
\]
Choose $c_0$ below the common lower bound of $\widetilde B$ and set
$\ell=(\widetilde B-c_0I)^{1/2}$. All square roots and inverses here are
uniformly analytic on a smaller strip. Finally put
\begin{align*}
 S'&=\mathsf A'(\widetilde B')^{-1}(\mathsf A')^*,\\
 \mathcal Mz&=\bigl(\mathsf A\widetilde B^{-1}T(\mathsf A')^*z,
                   a^{-1/2}d'^*(\mathsf A')^*z\bigr),\\
 \mathcal Nz&=\bigl((I-S')z,
 B_2^{1/2}\mathsf c'(\widetilde B')^{-1}(\mathsf A')^*z\bigr).
\end{align*}
Here $S'$ is an orthogonal projection. Substitution of
\eqref{free:Tidentities} proves \eqref{free:ambient} and
\eqref{free:harmonic}; $\mathsf c'd'=0$ proves \eqref{free:null}.
Using the definition of $U_*$ gives
$\mathcal M^*\mathcal M+\mathcal N^*\mathcal N=I$.
The extra factors are uniformly summable convolution kernels, so they
preserve the proved bounds for $T$. At depth zero keep $B=\widetilde
B=I$ instead of making the fine curl correction. This deletes a
nonnegative subtraction from $U_*$ and thus only improves its positivity.
All finite-period operators are obtained by folding these same kernels.
\end{proof}


